\section{A cyclotomic test case already determined by the manuscript}
\label{audit:cyclo}

The incoming alternating-tail report asks, as the first test in its
complementary-sector Question 7, for
$\Imag Z(2,1;i,-i)$, where
\[
 Z(2,1;i,-i)=\sum_{k>m\ge1}\frac{i^k(-i)^m}{k^2m}.
\]
The full two-resolvent complementary-sector generator remains open.
However, this particular coefficient is already an exact consequence of
two results in the pinned unified manuscript. Recording that implication
is an integration correction, not a newly discovered period evaluation.

\begin{proposition}[Explicit complementary coefficient]\label{audit:coefficient}
Put $\lambda_3=\Imag\Li_3((1+i)/2)$. Then
\begin{equation}\label{audit:coefficient-value}
 \boxed{\displaystyle
 \Imag Z(2,1;i,-i)
   =\lambda_3-\frac{\pi\log^2 2}{32}-\frac{\pi^3}{128}.}
\end{equation}
Equivalently,
\begin{equation}\label{audit:coefficient-integral}
 \int_0^1\frac{\log t\log(1-t)}{1+t^2}\dd t
 =\lambda_3-\frac{\pi\log^2 2}{32}-\frac{\pi^3}{128}.
\end{equation}
\end{proposition}
\begin{proof}
Use the one-direction convention
$\Li_{a,b}(z)=\sum_{k>m\ge1}z^k/(k^a m^b)$.
The inverse-color theorem in
\path{chapters/04-depth.tex}, equation
\texttt{inverse:eq:mixedb1}, gives
\[
 Z(2,1;z,z^{-1})
 =2\Li_3(z)+\Li_{2,1}(z)+\Li_{1,2}(z)-\zeta(2)\Li_1(z).
\]
The iterated-integral shuffle identity is
$\Li_1(z)\Li_2(z)=\Li_{1,2}(z)+2\Li_{2,1}(z)$.
Hence
\begin{equation}\label{audit:coefficient-reduction}
 Z(2,1;z,z^{-1})
 =2\Li_3(z)+\Li_1(z)\Li_2(z)-\Li_{2,1}(z)-\zeta(2)\Li_1(z).
\end{equation}
At $z=i$, put $L=\log2$ and $G=\beta(2)$. The standard values are
\[
 \Li_1(i)=-L/2+i\pi/4,\quad
 \Li_2(i)=-\pi^2/48+iG,\quad
 \Imag\Li_3(i)=\pi^3/32.
\]
The proved height-one specialization in
\path{chapters/04-shuffle-parity.tex} is
\[
 \Imag\Li_{2,1}(i)
 =-GL/2+\pi L^2/32-\lambda_3+3\pi^3/128.
\]
Substituting these identities into
\eqref{audit:coefficient-reduction} cancels the Catalan terms and yields
\eqref{audit:coefficient-value}. Both inherited reductions are proved in
the source, not inferred from a numerical fit \cite{ProveItPin}.

For the integral form, use the same manuscript's inverse-color integral
representation, or insert the elementary Mellin formula for $k^{-2}$
directly in the absolutely convergent double sum. It gives
\[
 Z(2,1;i,-i)
 =i\int_0^1\frac{\log t\log(1-t)}{1-it}\dd t.
\]
Taking the imaginary part proves \eqref{audit:coefficient-integral}.
The endpoint logarithms are integrable and $1-it$ does not vanish on the
integration interval, so no limiting branch choice is hidden here.
\end{proof}

The numerical value is
\[
 0.2806725061826611243960680536045272754284\ldots.
\]
This is a finite reduction into the explicitly stated classical
trilogarithm collection. It does not reduce $\lambda_3$ further and does
not prove its independence from any smaller set of constants. Nor does
it settle the full complementary generator or the distinct Gaussian
$S_6$ and revised $S_8$ conjectures.

The proposed status update is therefore precise: retain Question 7 for
the full complementary family, mark its displayed first coefficient as
already determined by the manuscript, and include
\eqref{audit:coefficient-value} as the missing cross-reference. The source
proofs of inverse-color transport and height-one reduction deserve the
credit.
